% AUTO-GENERATED by experiments/make_api_doc.py
\subsection*{\texttt{edgeguard.attestation.merkle}}
\begin{itemize}[leftmargin=*,itemsep=1pt]
\item \textbf{class \texttt{MerkleNode}} --- A single node in the Merkle tree.
\item \textbf{class \texttt{MerkleTree}} --- Binary Merkle tree built over a sequence of layer weight hashes.
\begin{itemize}[leftmargin=1.2em,itemsep=0pt]
\item \texttt{\small from\_\allowbreak{}dict(\allowbreak{}data: 'dict') -> "'MerkleTree'"} --- Reconstruct a MerkleTree from to\_dict() output.
\item \texttt{\small from\_\allowbreak{}hashes(\allowbreak{}layer\_\allowbreak{}hashes: 'Sequence[bytes]') -> "'MerkleTree'"} --- Build from pre-computed 32-byte leaf hashes (e.g. loaded from disk).
\item \texttt{\small from\_\allowbreak{}weights(\allowbreak{}layer\_\allowbreak{}weights: 'dict[str,\allowbreak{} np.ndarray]') -> "'MerkleTree'"} --- Hash each weight tensor and build the Merkle tree.
\item \texttt{\small get\_\allowbreak{}proof(\allowbreak{}self,\allowbreak{} leaf\_\allowbreak{}index: 'int') -> 'list[tuple[str,\allowbreak{} bytes]]'} --- Return the Merkle inclusion proof for a given leaf.
\item \texttt{\small to\_\allowbreak{}dict(\allowbreak{}self) -> 'dict'} --- Serialize the tree's leaf hashes to a JSON-compatible dict.
\item \texttt{\small verify\_\allowbreak{}and\_\allowbreak{}localize(\allowbreak{}self,\allowbreak{} fresh\_\allowbreak{}hashes: 'Sequence[bytes]') -> 'tuple[bool,\allowbreak{} list[int]]'} --- Compare fresh layer hashes against the stored tree.
\end{itemize}
\end{itemize}
\subsection*{\texttt{edgeguard.attestation.manager}}
\begin{itemize}[leftmargin=*,itemsep=1pt]
\item \textbf{class \texttt{AttestationResult}} --- Return value from MerkleAttestationManager.verify().
\begin{itemize}[leftmargin=1.2em,itemsep=0pt]
\item \texttt{\small summary(\allowbreak{}self) -> 'str'}
\end{itemize}
\item \textbf{class \texttt{MerkleAttestationManager}} --- Manages the full lifecycle of Merkle-tree-based weight attestation.
\begin{itemize}[leftmargin=1.2em,itemsep=0pt]
\item \texttt{\small build(\allowbreak{}self,\allowbreak{} layer\_\allowbreak{}weights: 'dict[str,\allowbreak{} np.ndarray]',\allowbreak{} private\_\allowbreak{}key=None) -> 'bytes'} --- Hash all layer weight tensors, build the Merkle tree, sign the root.
\item \texttt{\small load(\allowbreak{}path: 'str | Path') -> "'MerkleAttestationManager'"} --- Load an attestation previously saved with save().
\item \texttt{\small save(\allowbreak{}self,\allowbreak{} path: 'str | Path') -> 'None'} --- Persist the attestation state to a JSON file.
\item \texttt{\small verify(\allowbreak{}self,\allowbreak{} layer\_\allowbreak{}weights: 'dict[str,\allowbreak{} np.ndarray]') -> 'AttestationResult'} --- Verify the integrity of the model weights.
\end{itemize}
\end{itemize}
\subsection*{\texttt{edgeguard.detection.statistics}}
\begin{itemize}[leftmargin=*,itemsep=1pt]
\item \textbf{class \texttt{LayerStats}} --- Calibration statistics for one layer, computed over INT8 activations.
\begin{itemize}[leftmargin=1.2em,itemsep=0pt]
\item \texttt{\small from\_\allowbreak{}dict(\allowbreak{}data: 'dict') -> "'LayerStats'"}
\item \texttt{\small to\_\allowbreak{}dict(\allowbreak{}self) -> 'dict'}
\end{itemize}
\item \textbf{class \texttt{StreamingInt8Stats}} --- Online (streaming) computation of mean, variance, and covariance
\begin{itemize}[leftmargin=1.2em,itemsep=0pt]
\item \texttt{\small finalize(\allowbreak{}self) -> 'LayerStats'} --- Compute final statistics and return a LayerStats object.
\item \texttt{\small update(\allowbreak{}self,\allowbreak{} activation: 'np.ndarray') -> 'None'} --- Feed one INT8 activation sample (any shape; will be flattened).
\end{itemize}
\item \texttt{\small combined\_\allowbreak{}score(\allowbreak{}activation: 'np.ndarray',\allowbreak{} stats: 'LayerStats',\allowbreak{} alpha: 'float' = 0.5) -> 'dict'} --- Return both scores together with a weighted combination.
\item \texttt{\small mahalanobis\_\allowbreak{}distance(\allowbreak{}activation: 'np.ndarray',\allowbreak{} stats: 'LayerStats') -> 'float'} --- Mahalanobis distance from the calibration distribution.
\item \texttt{\small zscore\_\allowbreak{}max(\allowbreak{}activation: 'np.ndarray',\allowbreak{} stats: 'LayerStats') -> 'float'} --- Maximum absolute z-score across all channels.
\end{itemize}
\subsection*{\texttt{edgeguard.detection.calibration}}
\begin{itemize}[leftmargin=*,itemsep=1pt]
\item \textbf{class \texttt{Calibrator}} --- Streams clean activations through StreamingInt8Stats accumulators.
\begin{itemize}[leftmargin=1.2em,itemsep=0pt]
\item \texttt{\small feed\_\allowbreak{}activations(\allowbreak{}self,\allowbreak{} layer\_\allowbreak{}activations: 'dict[str,\allowbreak{} np.ndarray]') -> 'None'} --- Feed one sample's worth of per-layer INT8 activations.
\item \texttt{\small finalize(\allowbreak{}self) -> 'dict[str,\allowbreak{} LayerStats]'} --- Compute final statistics from all accumulated samples.
\item \texttt{\small load(\allowbreak{}path: 'str | Path') -> 'dict[str,\allowbreak{} LayerStats]'} --- Load calibration statistics previously saved with save().
\item \texttt{\small run(\allowbreak{}self,\allowbreak{} inputs: 'Iterable') -> 'None'} --- Run calibration over an iterable of raw inputs.
\item \texttt{\small save(\allowbreak{}stats: 'dict[str,\allowbreak{} LayerStats]',\allowbreak{} path: 'str | Path') -> 'None'} --- Save calibration statistics to a JSON file.
\end{itemize}
\end{itemize}
\subsection*{\texttt{edgeguard.detection.monitor}}
\begin{itemize}[leftmargin=*,itemsep=1pt]
\item \textbf{class \texttt{LayerScore}} --- Anomaly scores for a single layer on one inference call.
\item \textbf{class \texttt{LayerwiseActivationMonitor}} --- Runs per-layer activation anomaly detection using calibrated INT8 statistics.
\begin{itemize}[leftmargin=1.2em,itemsep=0pt]
\item \texttt{\small from\_\allowbreak{}calibration\_\allowbreak{}file(\allowbreak{}path: 'str | Path',\allowbreak{} **kwargs) -> "'LayerwiseActivationMonitor'"} --- Load calibration from disk and instantiate monitor.
\item \texttt{\small inspect(\allowbreak{}self,\allowbreak{} layer\_\allowbreak{}activations: 'dict[str,\allowbreak{} np.ndarray]') -> 'MonitorResult'} --- Score one inference pass.
\item \texttt{\small tune\_\allowbreak{}thresholds(\allowbreak{}self,\allowbreak{} clean\_\allowbreak{}results: 'Sequence[MonitorResult]',\allowbreak{} fpr\_\allowbreak{}target: 'float' = 0.001) -> 'tuple[float,\allowbreak{} float]'} --- Auto-tune zscore and Mahalanobis thresholds on clean validation results.
\end{itemize}
\item \textbf{class \texttt{MonitorResult}} --- Per-inference output from LayerwiseActivationMonitor.
\begin{itemize}[leftmargin=1.2em,itemsep=0pt]
\item \texttt{\small summary(\allowbreak{}self) -> 'str'}
\end{itemize}
\end{itemize}
\subsection*{\texttt{edgeguard.classifier.classifier}}
\begin{itemize}[leftmargin=*,itemsep=1pt]
\item \textbf{class \texttt{ThreatClassifier}} --- Fuses Merkle attestation results and activation monitor results into
\begin{itemize}[leftmargin=1.2em,itemsep=0pt]
\item \texttt{\small classify(\allowbreak{}self,\allowbreak{} monitor\_\allowbreak{}result: 'MonitorResult',\allowbreak{} attestation\_\allowbreak{}result: 'Optional[AttestationResult]' = None) -> 'Threa...)} --- Classify the threat and generate a layerwise explanation.
\end{itemize}
\item \textbf{class \texttt{ThreatResult}} --- Full output of the ThreatClassifier.
\begin{itemize}[leftmargin=1.2em,itemsep=0pt]
\item \texttt{\small is\_\allowbreak{}threat(\allowbreak{}self) -> 'bool'}
\item \texttt{\small summary(\allowbreak{}self) -> 'str'}
\end{itemize}
\item \textbf{class \texttt{ThreatType}} --- Create a collection of name/value pairs.
\end{itemize}
\subsection*{\texttt{edgeguard.classifier.rules}}
\begin{itemize}[leftmargin=*,itemsep=1pt]
\item \textbf{class \texttt{ThreatRules}} --- Thresholds and scoring weights for threat classification.
\begin{itemize}[leftmargin=1.2em,itemsep=0pt]
\item \texttt{\small adversarial\_\allowbreak{}confidence(\allowbreak{}self,\allowbreak{} early\_\allowbreak{}ratio: 'float',\allowbreak{} max\_\allowbreak{}zscore: 'float') -> 'float'} --- Confidence that the threat is ADVERSARIAL\_INPUT.
\item \texttt{\small permissive(\allowbreak{}) -> "'ThreatRules'"} --- Rules tuned for low false-positive rate (miss more subtle threats).
\item \texttt{\small strict(\allowbreak{}) -> "'ThreatRules'"} --- Rules tuned for high recall (may have more false positives).
\item \texttt{\small transient\_\allowbreak{}fault\_\allowbreak{}confidence(\allowbreak{}self,\allowbreak{} early\_\allowbreak{}ratio: 'float',\allowbreak{} max\_\allowbreak{}zscore: 'float') -> 'float'} --- Confidence that the threat is TRANSIENT\_FAULT.
\item \texttt{\small weight\_\allowbreak{}tamper\_\allowbreak{}confidence(\allowbreak{}self,\allowbreak{} n\_\allowbreak{}tampered: 'int',\allowbreak{} n\_\allowbreak{}total: 'int',\allowbreak{} act\_\allowbreak{}anomaly: 'bool') -> 'float'} --- Confidence that the threat is WEIGHT\_TAMPER.
\end{itemize}
\end{itemize}
\subsection*{\texttt{edgeguard.llm.model}}
\begin{itemize}[leftmargin=*,itemsep=1pt]
\item \textbf{class \texttt{LLMWeightProfile}} --- Describes the architecture of a quantized transformer LLM.
\item \textbf{class \texttt{QuantizedTransformerStub}} --- Generates a realistic INT8 weight dictionary mirroring a transformer LLM.
\begin{itemize}[leftmargin=1.2em,itemsep=0pt]
\item \texttt{\small get\_\allowbreak{}weights(\allowbreak{}self) -> 'dict[str,\allowbreak{} np.ndarray]'} --- Generate (or return cached) INT8 weight tensors.
\end{itemize}
\end{itemize}
\subsection*{\texttt{edgeguard.llm.detector}}
\begin{itemize}[leftmargin=*,itemsep=1pt]
\item \textbf{class \texttt{LLMThreatDetector}} --- Unified threat detector for quantized LLMs on edge devices.
\begin{itemize}[leftmargin=1.2em,itemsep=0pt]
\item \texttt{\small detect(\allowbreak{}self,\allowbreak{} activations: 'dict[str,\allowbreak{} np.ndarray]',\allowbreak{} weights: 'Optional[dict[str,\allowbreak{} np.ndarray]]' = None) -> 'LLMThre...)} --- Run the full LLM threat detection pipeline.
\end{itemize}
\item \textbf{class \texttt{LLMThreatResult}} --- Extended threat result for LLM-specific threats.
\begin{itemize}[leftmargin=1.2em,itemsep=0pt]
\item \texttt{\small is\_\allowbreak{}threat(\allowbreak{}self) -> 'bool'}
\item \texttt{\small summary(\allowbreak{}self) -> 'str'}
\end{itemize}
\item \textbf{class \texttt{LLMThreatType}} --- Create a collection of name/value pairs.
\end{itemize}
\subsection*{\texttt{edgeguard.attacks.bitflip}}
\begin{itemize}[leftmargin=*,itemsep=1pt]
\item \textbf{class \texttt{BitFlipAttack}} --- Simulates Rowhammer-style bit-flip attacks on INT8 quantized model weights.
\begin{itemize}[leftmargin=1.2em,itemsep=0pt]
\item \texttt{\small attack(\allowbreak{}self,\allowbreak{} weights: 'dict[str,\allowbreak{} np.ndarray]',\allowbreak{} target\_\allowbreak{}layer: 'Optional[str]' = None) -> 'BitFlipResult'} --- Apply bit-flip attack to the weight dictionary.
\item \texttt{\small weight\_\allowbreak{}delta\_\allowbreak{}stats(\allowbreak{}original: 'dict[str,\allowbreak{} np.ndarray]',\allowbreak{} attacked: 'dict[str,\allowbreak{} np.ndarray]') -> 'dict'} --- Compute per-layer weight change statistics between original and attacked.
\end{itemize}
\item \textbf{class \texttt{BitFlipResult}} --- Result of a bit-flip attack simulation.
\begin{itemize}[leftmargin=1.2em,itemsep=0pt]
\item \texttt{\small magnitude\_\allowbreak{}shift(\allowbreak{}self) -> 'dict[str,\allowbreak{} float]'} --- Not computable without original --- see BitFlipAttack.attack().
\end{itemize}
\end{itemize}
\subsection*{\texttt{edgeguard.attacks.weight\_extraction}}
\begin{itemize}[leftmargin=*,itemsep=1pt]
\item \textbf{class \texttt{ExtractionResult}} --- Result of a weight extraction simulation.
\begin{itemize}[leftmargin=1.2em,itemsep=0pt]
\item \texttt{\small summary(\allowbreak{}self) -> 'str'}
\end{itemize}
\item \textbf{class \texttt{WeightExtractionAttack}} --- Simulates read-based weight extraction attacks.
\begin{itemize}[leftmargin=1.2em,itemsep=0pt]
\item \texttt{\small extract(\allowbreak{}self,\allowbreak{} weights: 'dict[str,\allowbreak{} np.ndarray]',\allowbreak{} quant\_\allowbreak{}scales: 'Optional[dict[str,\allowbreak{} float]]' = None,\allowbreak{} quant\_\allowbreak{}zero\_\allowbreak{}poin...)} --- Simulate extraction of model weights.
\end{itemize}
\end{itemize}
\subsection*{\texttt{edgeguard.attacks.training\_data\_recovery}}
\begin{itemize}[leftmargin=*,itemsep=1pt]
\item \textbf{class \texttt{RecoveryResult}} --- Simulation of a training data recovery attempt.
\begin{itemize}[leftmargin=1.2em,itemsep=0pt]
\item \texttt{\small summary(\allowbreak{}self) -> 'str'}
\end{itemize}
\item \textbf{class \texttt{TrainingDataRecovery}} --- Simulates training-data extraction probing and generates realistic
\begin{itemize}[leftmargin=1.2em,itemsep=0pt]
\item \texttt{\small compare\_\allowbreak{}signatures(\allowbreak{}normal: 'dict[str,\allowbreak{} np.ndarray]',\allowbreak{} extraction: 'dict[str,\allowbreak{} np.ndarray]') -> 'dict'} --- Compare activation signatures between normal and extraction-probe passes.
\item \texttt{\small normal\_\allowbreak{}inference\_\allowbreak{}activation(\allowbreak{}self,\allowbreak{} layer\_\allowbreak{}names: 'list[str]',\allowbreak{} activation\_\allowbreak{}dim: 'int' = 64) -> 'dict[str,\allowbreak{} np.ndarray]'} --- Generate activations for a normal (non-extracting) inference pass.
\item \texttt{\small simulate\_\allowbreak{}probe(\allowbreak{}self,\allowbreak{} layer\_\allowbreak{}names: 'list[str]',\allowbreak{} n\_\allowbreak{}queries: 'int' = 100,\allowbreak{} activation\_\allowbreak{}dim: 'int' = 64) -> 'RecoveryResult'} --- Simulate a training-data extraction attack and return the resulting
\end{itemize}
\end{itemize}
\subsection*{\texttt{edgeguard.attacks.supply\_chain}}
\begin{itemize}[leftmargin=*,itemsep=1pt]
\item \textbf{class \texttt{PoisoningResult}} --- Result of a supply-chain poisoning attack.
\begin{itemize}[leftmargin=1.2em,itemsep=0pt]
\item \texttt{\small summary(\allowbreak{}self) -> 'str'}
\end{itemize}
\item \textbf{class \texttt{SupplyChainPoisoning}} --- Simulates supply-chain weight poisoning attacks.
\begin{itemize}[leftmargin=1.2em,itemsep=0pt]
\item \texttt{\small poison(\allowbreak{}self,\allowbreak{} weights: 'dict[str,\allowbreak{} np.ndarray]',\allowbreak{} target\_\allowbreak{}layer: 'Optional[str]' = None) -> 'PoisoningResult'} --- Apply supply-chain poisoning to weight tensors.
\end{itemize}
\end{itemize}
\subsection*{\texttt{edgeguard.attacks.adversarial}}
\begin{itemize}[leftmargin=*,itemsep=1pt]
\item \textbf{class \texttt{AdversarialInputAttack}} --- Simulates adversarial input attacks and their distinctive activation patterns.
\begin{itemize}[leftmargin=1.2em,itemsep=0pt]
\item \texttt{\small attack(\allowbreak{}self,\allowbreak{} layer\_\allowbreak{}names: 'list[str]',\allowbreak{} activation\_\allowbreak{}dim: 'int' = 64) -> 'AdversarialResult'} --- Simulate an adversarial attack and generate characteristic activations.
\end{itemize}
\item \textbf{class \texttt{AdversarialResult}} --- Result of an adversarial input attack simulation.
\begin{itemize}[leftmargin=1.2em,itemsep=0pt]
\item \texttt{\small summary(\allowbreak{}self) -> 'str'}
\end{itemize}
\end{itemize}
\subsection*{\texttt{edgeguard.attacks.synergy}}
\begin{itemize}[leftmargin=*,itemsep=1pt]
\item \textbf{class \texttt{SynergyAttack}} --- Simulates the combined BFA + adversarial input attack.
\begin{itemize}[leftmargin=1.2em,itemsep=0pt]
\item \texttt{\small attack(\allowbreak{}self,\allowbreak{} weights: 'dict[str,\allowbreak{} np.ndarray]',\allowbreak{} layer\_\allowbreak{}names: 'list[str]',\allowbreak{} target\_\allowbreak{}layer: 'Optional[str]' = None,\allowbreak{} ac...)} --- Execute the combined attack.
\end{itemize}
\item \textbf{class \texttt{SynergyResult}} --- Result of a combined BFA + adversarial attack.
\begin{itemize}[leftmargin=1.2em,itemsep=0pt]
\item \texttt{\small summary(\allowbreak{}self) -> 'str'}
\end{itemize}
\end{itemize}
\subsection*{\texttt{edgeguard.instrumentation.hooks}}
\begin{itemize}[leftmargin=*,itemsep=1pt]
\item \textbf{class \texttt{ActivationCapture}} --- A callable that can be registered as a PyTorch forward hook.
\begin{itemize}[leftmargin=1.2em,itemsep=0pt]
\item \texttt{\small reset(\allowbreak{}self) -> 'None'} --- Clear the buffered activation.
\end{itemize}
\item \textbf{class \texttt{HookManager}} --- Manages a set of LayerHooks and collects activations after each forward pass.
\begin{itemize}[leftmargin=1.2em,itemsep=0pt]
\item \texttt{\small collect(\allowbreak{}self) -> 'dict[str,\allowbreak{} np.ndarray]'} --- Return the most recent activations from all registered hooks.
\item \texttt{\small register(\allowbreak{}self,\allowbreak{} layer\_\allowbreak{}map: 'dict[str,\allowbreak{} object]') -> "'HookManager'"} --- Register hooks for each named layer.
\item \texttt{\small remove\_\allowbreak{}all(\allowbreak{}self) -> 'None'} --- Remove all registered hooks.
\item \texttt{\small reset(\allowbreak{}self) -> 'None'} --- Clear all buffered activations.
\end{itemize}
\item \textbf{class \texttt{LayerHook}} --- Wraps an ActivationCapture and its PyTorch hook handle together.
\begin{itemize}[leftmargin=1.2em,itemsep=0pt]
\item \texttt{\small register(\allowbreak{}self,\allowbreak{} module) -> "'LayerHook'"} --- Register the hook on a PyTorch module.
\item \texttt{\small remove(\allowbreak{}self) -> 'None'} --- Remove the hook from the module.
\end{itemize}
\end{itemize}
\subsection*{\texttt{edgeguard.instrumentation.qnn\_adapter}}
\begin{itemize}[leftmargin=*,itemsep=1pt]
\item \textbf{class \texttt{InstrumentedQNNRunner}} --- Wraps a QNNGraph with activation capture for use with EdgeGuard.
\begin{itemize}[leftmargin=1.2em,itemsep=0pt]
\item \texttt{\small run(\allowbreak{}self,\allowbreak{} inputs: 'Dict[str,\allowbreak{} np.ndarray]') -> 'Dict[str,\allowbreak{} np.ndarray]'} --- Execute one inference and return per-layer INT8 activations.
\end{itemize}
\item \textbf{class \texttt{QNNActivationBuffer}} --- Thin ring buffer that stores per-layer INT8 activations from QNN inference.
\begin{itemize}[leftmargin=1.2em,itemsep=0pt]
\item \texttt{\small capture(\allowbreak{}self,\allowbreak{} graph: 'QNNGraph') -> 'None'} --- Pull the latest output tensors from the QNN graph into the buffer.
\item \texttt{\small latest(\allowbreak{}self) -> 'Dict[str,\allowbreak{} np.ndarray]'} --- Return the most recent activation for each layer.
\item \texttt{\small reset(\allowbreak{}self) -> 'None'} --- Clear all buffered activations.
\end{itemize}
\item \textbf{class \texttt{QNNAdapterConfig}} --- Configuration for QNN instrumentation.
\item \textbf{class \texttt{QNNGraph}} --- Minimal interface to a compiled QNN graph (duck-typed).
\begin{itemize}[leftmargin=1.2em,itemsep=0pt]
\item \texttt{\small execute(\allowbreak{}self,\allowbreak{} inputs: 'Dict[str,\allowbreak{} np.ndarray]') -> 'Dict[str,\allowbreak{} np.ndarray]'} --- Run one inference pass.
\item \texttt{\small get\_\allowbreak{}op\_\allowbreak{}names(\allowbreak{}self) -> 'List[str]'} --- Return all op node names in the graph.
\item \texttt{\small get\_\allowbreak{}output\_\allowbreak{}tensor(\allowbreak{}self,\allowbreak{} op\_\allowbreak{}name: 'str') -> 'Optional[np.ndarray]'} --- Return the most recent output tensor for a given op (after execute()).
\end{itemize}
\item \texttt{\small build\_\allowbreak{}layer\_\allowbreak{}map(\allowbreak{}graph: 'QNNGraph',\allowbreak{} target\_\allowbreak{}patterns: 'Sequence[str]') -> 'Dict[str,\allowbreak{} List[str]]'} --- Auto-build a layer\_op\_map by matching QNN op names against patterns.
\end{itemize}
\subsection*{\texttt{edgeguard.utils.crypto}}
\begin{itemize}[leftmargin=*,itemsep=1pt]
\item \texttt{\small generate\_\allowbreak{}keypair(\allowbreak{}) -> 'tuple[ec.EllipticCurvePrivateKey,\allowbreak{} ec.EllipticCurvePublicKey]'} --- Generate a fresh P-256 keypair for signing Merkle roots.
\item \texttt{\small hash\_\allowbreak{}pair(\allowbreak{}left: 'bytes',\allowbreak{} right: 'bytes') -> 'bytes'} --- Combine two 32-byte digests into a parent node hash.
\item \texttt{\small hash\_\allowbreak{}weight\_\allowbreak{}tensor(\allowbreak{}weights: 'np.ndarray') -> 'bytes'} --- Compute a SHA-256 digest over the raw bytes of a quantized weight tensor.
\item \texttt{\small private\_\allowbreak{}key\_\allowbreak{}from\_\allowbreak{}pem(\allowbreak{}pem: 'bytes') -> 'ec.EllipticCurvePrivateKey'} --- 
\item \texttt{\small private\_\allowbreak{}key\_\allowbreak{}to\_\allowbreak{}pem(\allowbreak{}private\_\allowbreak{}key: 'ec.EllipticCurvePrivateKey') -> 'bytes'} --- 
\item \texttt{\small public\_\allowbreak{}key\_\allowbreak{}from\_\allowbreak{}pem(\allowbreak{}pem: 'bytes') -> 'ec.EllipticCurvePublicKey'} --- 
\item \texttt{\small public\_\allowbreak{}key\_\allowbreak{}to\_\allowbreak{}pem(\allowbreak{}public\_\allowbreak{}key: 'ec.EllipticCurvePublicKey') -> 'bytes'} --- 
\item \texttt{\small sign\_\allowbreak{}root(\allowbreak{}root\_\allowbreak{}hash: 'bytes',\allowbreak{} private\_\allowbreak{}key: 'ec.EllipticCurvePrivateKey') -> 'bytes'} --- ECDSA-SHA256 signature over a Merkle root hash.
\item \texttt{\small verify\_\allowbreak{}root(\allowbreak{}root\_\allowbreak{}hash: 'bytes',\allowbreak{} signature: 'bytes',\allowbreak{} public\_\allowbreak{}key: 'ec.EllipticCurvePublicKey') -> 'bool'} --- Verify an ECDSA-SHA256 signature over a Merkle root hash.
\end{itemize}
\subsection*{\texttt{edgeguard.utils.quantization}}
\begin{itemize}[leftmargin=*,itemsep=1pt]
\item \textbf{class \texttt{QuantizationParams}} --- Scale and zero-point for affine quantization: q = round(x / scale) + zero\_point.
\item \texttt{\small compute\_\allowbreak{}int8\_\allowbreak{}params(\allowbreak{}tensor\_\allowbreak{}fp32: 'np.ndarray',\allowbreak{} symmetric: 'bool' = True) -> 'QuantizationParams'} --- Derive per-tensor INT8 quantization parameters from a floating-point tensor.
\item \texttt{\small dequantize(\allowbreak{}tensor\_\allowbreak{}int: 'np.ndarray',\allowbreak{} params: 'QuantizationParams') -> 'np.ndarray'} --- Reverse affine quantization: x $\approx$ (q - zero\_point) * scale.
\item \texttt{\small int4\_\allowbreak{}pack(\allowbreak{}weights\_\allowbreak{}int8: 'np.ndarray') -> 'np.ndarray'} --- Pack an array of INT4 values stored in int8 (range -8..7) into nibbles.
\item \texttt{\small int4\_\allowbreak{}unpack(\allowbreak{}packed: 'np.ndarray',\allowbreak{} length: 'int') -> 'np.ndarray'} --- Unpack nibbles back to int8 values (sign-extended from 4 bits).
\item \texttt{\small layer\_\allowbreak{}weight\_\allowbreak{}stats(\allowbreak{}weights\_\allowbreak{}int8: 'np.ndarray') -> 'dict'} --- Compute basic statistics over quantized layer weights.
\item \texttt{\small quantize(\allowbreak{}tensor\_\allowbreak{}fp32: 'np.ndarray',\allowbreak{} params: 'QuantizationParams') -> 'np.ndarray'} --- Affine-quantize a float32 tensor to int8/uint8.
\end{itemize}
